%!TEX root = ../../../main.tex

\chapter{Architektur}
\label{chap:architektur}

\section{Überblick}

Abbildung~\ref{fig:architektur} zeigt den Aufbau: Die Dokumente werden einmal eingelesen und in drei Ebenen gespeichert. Zwei Prototypen beantworten die Fragen mit derselben Suche und denselben Kerndaten, P1 mit einem festen Ablauf, den Code steuert, P2 als ReAct-Agent. Die Modelle laufen in Ollama \cite{ollama}, Vektorspeicher ist ChromaDB \cite{chromadb}, die Oberfläche eine Streamlit-App.

\begin{figure}[tb]
\centering
\begin{tikzpicture}[font=\scriptsize, node distance=3mm and 3.5mm,
    box/.style={draw, rounded corners=2pt, align=center, minimum height=8mm, inner sep=2pt},
    store/.style={box, fill=gray!12}, llmbox/.style={box, fill=blue!8}, codebox/.style={box, fill=green!8},
    arr/.style={-{Stealth[length=1.6mm]}}]
  \node[box] (raw) {PDF / DOCX\\\texttt{data/raw}};
  \node[box, right=of raw] (parse) {Parser\\Tabellenzeilen,\\Kapitelpfad, Bilder};
  \node[store, right=of parse] (json) {Ebene 1: JSON\\je Dokument};
  \node[llmbox, below=of json] (enrich) {Bildbeschreibung (gemma3)\\Kerndaten (qwen2.5)};
  \node[store, left=of enrich] (cache) {Ebene 2: Cache\\SHA-256};
  \node[store, right=of json] (chroma) {Ebene 3: ChromaDB\\nomic-embed-text};
  \node[codebox, below=of chroma] (cat) {Katalog\\16 Dokumente,\\Kerndaten};
  % P1
  \node[llmbox, below=9mm of cache, xshift=-4mm] (plan) {1 Plan\\(qwen2.5)};
  \node[codebox, right=of plan] (bestand) {2 Bestand};
  \node[codebox, right=of bestand] (tool) {3 Werkzeug\\Suche / Filter};
  \node[codebox, right=of tool] (bel) {4 Belege};
  \node[llmbox, right=of bel] (draft) {5 Entwurf\\(qwen2.5)};
  \node[codebox, right=of draft] (check) {6 Prüfung};
  \node[left=1mm of plan, font=\scriptsize\bfseries] {P1};
  \node[llmbox, below=4mm of tool, minimum width=6cm] (p2) {P2: ReAct (llama3.1) wählt aus 8 Werkzeugen};
  \node[left=1mm of p2, font=\scriptsize\bfseries] {P2};
  \draw[arr] (raw) -- (parse); \draw[arr] (parse) -- (json); \draw[arr] (json) -- (enrich);
  \draw[arr, <->] (enrich) -- (cache); \draw[arr] (json) -- (chroma); \draw[arr] (enrich) -- (cat);
  \draw[arr] (plan) -- (bestand); \draw[arr] (bestand) -- (tool); \draw[arr] (tool) -- (bel);
  \draw[arr, dashed] (bel) -- node[above, font=\tiny]{Klick} (draft); \draw[arr] (draft) -- (check);
  \draw[arr, densely dotted] (chroma.south east) |- ++(0,-4mm) -| (tool.north);
  \draw[arr, densely dotted] (cat.south) |- (plan.north);
\end{tikzpicture}
\caption{Architektur: Einlesen in drei Speicherebenen, darunter die sechs Stufen von P1 (grün: Code, blau: Sprachmodell) und P2}
\label{fig:architektur}
\end{figure}

\section{Vorverarbeitung und Speicherebenen}

Der Parser liest PDF mit PyMuPDF und DOCX mit python-docx. Zuerst sucht er Tabellen und schneidet sie aus dem Fließtext aus. Jede Tabellenzeile wird ein eigener Block wie \enquote{Gebührenart: Schmutzwassergebühr | Menge: 125,0\,m³ | Satz: 2,80\,€/m³ | Brutto (€): 416,50}. Ein normaler Text-Splitter würde die Zeile zerschneiden, und das Modell wüsste nicht mehr, ob eine Zahl netto oder brutto ist. Überschriften erkennt der Parser an Schriftgröße oder Fettdruck und gibt jedem Block seinen Kapitelpfad mit. Absätze unter 20 Zeichen hängt er an den vorigen an, damit ein Betrag nicht allein steht, über 1\,500 Zeichen teilt er sie. gemma3:4b beschreibt Bilder mit allen ablesbaren Werten und fasst jede Tabelle zusammen. Aus 16 Dokumenten entstehen 436 Blöcke.

Ebene~1 ist eine lesbare JSON-Datei je Dokument. Dort sehe ich nach, was der Parser gemacht hat. Ebene~2 ist ein Cache für Modellausgaben mit dem SHA-256-Hash des Blocks als Schlüssel. Ändere ich den Parser, beschreibt das Modell nur geänderte Blöcke neu. Ebene~3 ist ChromaDB mit Embeddings von nomic-embed-text \cite{nussbaum2024}. Zusätzlich liest qwen2.5:7b je Dokument 19 \emph{Kerndaten} aus, etwa Anbieter, Kunde, Beträge und Fristen. Steht ein Wert nicht wörtlich im Dokument, löscht Code ihn. Danach stimmen 96\,\% der Felder. Daraus baut Code einen \emph{Katalog} mit einer Zeile je Dokument, etwa \enquote{[12] Vertrag · Volt \& Wärme AG · Heizstromvertrag WP-Natur}.

\section{Suche}

Beim Einbetten steht vor jedem Block sein Pfad \texttt{[Ordner | Datei | Kapitel]}. So wird \enquote{2,80\,€/m³} als Teil der \emph{Wasserabrechnung} gespeichert und nicht als irgendeine Zahl. ChromaDB liefert 20 Kandidaten, Treffer, deren Pfad ein Wort der Frage enthält, zählen 1,5-fach. Drei Zusätze aus dem modularen RAG habe ich gemessen und wieder entfernt (Anhang~\ref{app:messungen}): BM25 \cite{robertson2009} brachte nichts, weil der Pfad die Fachwörter schon enthält, ein Cross-Encoder \cite{nogueira2019} machte die Reihenfolge schlechter, und das Umformulieren der Frage senkte den Recall von 100 auf 68\,\%.

\section{Prototyp 1: Orchestrator mit Code-Prüfung}
\label{sec:p1}

P1 arbeitet in sechs festen Stufen (Abbildung~\ref{fig:architektur}), das Modell kommt nur zweimal vor.

\begin{enumerate}[nosep, leftmargin=*]
    \item \emph{Plan:} qwen2.5:7b sieht den Katalog und füllt ein kleines JSON: Fragetyp (\texttt{wert}, \texttt{vergleich}, \texttt{frist}, \texttt{liste}, \texttt{fehlt}), Dokumente und bei Listen einen Filter. Code prüft den Plan, ergänzt etwa ein Dokument, das die Frage eindeutig nennt, und verlangt bei Vergleichen Rechnung \emph{und} Vertrag.
    \item \emph{Bestand:} Code prüft, ob es die gefragte Dokumentart gibt. Bei \enquote{Gasvertrag} antwortet P1 sofort \enquote{nicht belegt}.
    \item \emph{Werkzeug:} Werte sucht P1 nur in den geplanten Dokumenten. Listen filtert Python über die Kerndaten, Kündigungstermine berechnet ein Fristrechner.
    \item \emph{Belege:} Die App zeigt die nummerierten Treffer zuerst, nach Ordner, Datei und Kapitel geordnet und nach Relevanz eingefärbt (grün ab 0,7, rot unter 0,4). Erst dann klicke ich auf \enquote{Antwort}.
    \item \emph{Entwurf:} qwen2.5:7b liefert als JSON die Werte mit Belegnummer, optional eine Rechnung (\enquote{w1 $-$ w2}) und einen Antwortsatz.
    \item \emph{Prüfung:} Code sucht jeden Wert im genannten Beleg, prüft die Spalte (netto oder brutto) und rechnet selbst. Fehler gehen einmal konkret an das Modell zurück (\enquote{14,00 steht nicht in Beleg~1}); was dann noch unbelegt ist, fällt weg.
\end{enumerate}

Das Modell soll also nur lesen. Rechnen und entscheiden, ob etwas belegt ist, macht Code. Das ähnelt Self-RAG \cite{asai2024}, nur prüft bei mir nicht das Modell sich selbst, denn ein 7B-Modell erkennt eigene Fehler schlecht.

\section{Prototyp 2: ReAct-Agent}

P2 folgt dem ReAct-Muster \cite{yao2023}: llama3.1:8b überlegt, ruft ein Werkzeug auf, liest das Ergebnis und entscheidet, ob es noch eins braucht. Die acht Werkzeuge sind Suche in Text, Tabellen und Bildbeschreibungen, Gliederung eines Dokuments, Rechner (nach der Idee von Toolformer \cite{schick2023}), Kerndaten, Vertrags- und Rechnungsliste. Zahlprüfung und Quellenliste sind dieselben wie bei P1.

\section{Weg einer Frage durch P1}
\label{sec:weg}

Testfrage T14 lautet \enquote{Um wie viel Euro ist die neue quartalsweise Wasser-Vorauszahlung ab Februar 2025 höher als die Vorauszahlung im Wasservertrag?} Der \emph{Plan} wählt in 17\,s Wasserrechnung und -vertrag. Der Katalog liegt im Prompt-Cache, das Modell liest nur 23 neue Tokens. \emph{Bestand} und \emph{Suche} liefern in 0,2\,s zwölf Belege, darunter 210,00\,€ aus §\,5 des Vertrags und \enquote{Ab dem 15.02.2025 \dots{} auf 225,00\,€ festgesetzt} aus der Rechnung. Der \emph{Entwurf} braucht 113\,s, davon 79\,s nur zum Lesen von rund 1\,500 Tokens Belegen. Er nennt beide Werte mit richtiger Belegnummer und antwortet 15,00\,€, die \emph{Prüfung} findet beide Werte. Aber: Das Modell hat die Differenz selbst ausgerechnet, statt den Platzhalter für das Code-Ergebnis zu nutzen. Die Zahlprüfung ließ \enquote{15} nur durch, weil die Zahl als Tag im Datum 15.02.2025 im Beleg steht. Die Antwort ist richtig, aber aus dem falschen Grund geprüft.

\section{Warum kein GraphRAG}
\label{sec:graphrag}

\emph{Entitäten und Beziehungen:} Ein Graph hätte Personen, Anbieter, Verträge, Rechnungen, Zähler und Fahrzeuge als Knoten und Kanten wie \enquote{Rechnung gehört zu Vertrag} oder \enquote{Vertrag läuft auf Person}. \emph{Profitieren würden} vor allem Typ c und d, etwa \enquote{Welche Verträge laufen auf Max Mustermann?}, und Fragen über mehrere Schritte. \emph{Klassisches RAG reicht trotzdem:} Bei 16 Dokumenten ist jede Kante ein Vergleich in den Kerndaten (gleiche Vertragsnummer, Feld \enquote{Kunde}), den Plan und Filter erledigen. Themenfragen über die ganze Sammlung, für die GraphRAG gedacht ist, kommen bei mir nicht vor. \emph{Kosten:} GraphRAG braucht beim Einlesen einen Modellaufruf je Block, bei 436 Blöcken und geschätzt 30\,s je Aufruf über drei Stunden; die Kerndaten brauchen 16 Aufrufe. Meine erste Version hatte fünf Graph-Werkzeuge. Keine Testfrage brauchte sie, also habe ich sie entfernt.
